%\chapter{Background equations and formulae}

\section{Background equations and formulae}

\subsection{Coordinate systems, surface averages and fluxes}

In this section we provide a summary of the main relations used in this
paper. Detailed derivations and discussions can be found in \cite{HH}. 
For the purpose of the geometry description, we consider the magnetic
system as being toroidally symmetric. 
Small violations of symmetry, which are always present in real systems,
are neglected in the equilibrium, but may taken into account in the confinement
properties of the plasma. 

For the description of tokamak geometry we use two coordinate systems.
The first one is the cylindrical coordinate system $\{r,\vf,z\}$ with
the polar axis coinciding with the major axis of a torus.  
The second $\{a,\theta,\zeta\}$ is related to the magnetic geometry of
the tokamak, and uses a ``radial'' variable $a,$ which is an arbitrary
label of a magnetic flux surface, and a poloidal angle $\theta,$ the
specification of which does not play any role below. 
In order to make both systems right-handed we choose the toroidal angle
$\zeta$ as $\zeta=-\vf.$ 

We denote a flux surface, defined by the equation $a={\rm Const}$, as
$\mbox{\msym S}_a$ and introduce a volume integral over the interior,
{\msym V}, of this flux surface 
\begin{equation}
\label{dfV}
V=\int\limits_{\mbox{\msym V}} dV
=\int\limits_{0}^{a}\! da \int\limits_{\mbox{\msym S}_a}
\frac{dS_a}{|\nb a|} 
=\int\limits_{0}^{a}\! da \int\limits_{0}^{2\pi}\!
d\zeta\int\limits_{0}^{2\pi}\sqrt{g} \, d\theta 
=2\pi\int\limits_{0}^{a}\! da\int\limits_{0}^{2\pi}\sqrt{g} \, d\theta
\end{equation}
where $g$ is the determinant of the metric tensor 
$g
=\dstyle\left(\frac{D(x,y,z)}{D(a,\theta,\zeta)}\right)^2
=\left(\nb a\nb\theta\nb\zeta\right)^{-2}
$ and $V$ is the volume of {\msym V}. 
The flux surface average of an arbitrary function $f(\br)$ is defined
\begin{equation}
\label{dfav}
\av{f}
=\prti{}{V}\int\limits_{\mbox{\msym V}} f\, dV
=\prti{}{V}\int\limits_{0}^{a}\! da \int\limits_{0}^{2\pi}\!
d\zeta\int\limits_{0}^{2\pi}\sqrt{g} f\, d\theta 
=2\pi\prti{a}{V}\int\limits_{0}^{2\pi}\sqrt{g} f\, d\theta
.\end{equation}
The symbols for partial derivatives are used here in order to emphasize
that all surface functions (i.e. the functions of single space
coordinate $a$) are also functions of time $t$.  
Equation (\ref{dfav}) can equivalently be represented as
\begin{equation}
\label{dfav2}
\av{f}
=\prti{}{V}\int\limits_{\mbox{\msym V}} f\, dV
=\int\limits_{\mbox{\msym S}_a}f\frac{dS_a}{|\nb V|}
=\prti{\psi}{V}\oint f\frac{dl_\theta}{B_{\rm pol}}
=\left.\oint f\frac{dl_\theta}{B_{\rm pol}}\right/
\oint \frac{dl_\theta}{B_{\rm pol}}
,\end{equation}
where the poloidal flux $\psi$ and the poloidal component, $B_{\rm
pol},$ of the magnetic field {\bf B} are determined in the next
section. 
For any axisymmetric function, $f=f(a,\theta),$ one also has
\begin{equation}
\label{avBf}
\av{{\bf B}\.\na f}
=\av{\mbox{div}\lb(f{\bf B}\rb)}
=\prti{}{V}\int\limits_{\mbox{\msym V}}\mbox{div}\lb(f{\bf B}\rb)dV
=\prti{}{V}\oint\limits_{\mbox{\msym S}_a}f\,{\bf B}\.d{\bf s}
=0
.\end{equation}
Making use of \req{dfav2} we find the area of a magnetic surface
\begin{equation}\label{Sa}
S_a
=\oint\limits_{\mbox{\msym S}_a}dS_a
=\prti{}{V}\int\limits_{\mbox{\msym V}}|\na V|dV
=\av{|\na V|}
=\prti{V}{a}\av{|\na a|}
.\end{equation}
Further useful properties of the averaging procedure are
\begin{equation}\label{avdiv}
\av{F(a)f(a,\theta)}=F(a)\av{f}
,\qquad
\prti{V}{a}=2\pi\int\limits_{0}^{2\pi}\sqrt{g}\,d\theta
,\qquad
\av{\mbox{div}\,{\bf g}}
=\prti{}{V}\av{{\bf g}\.\na V}
=\prti{\mit\Gamma}{V}
.\end{equation}
The last identity in \req{avdiv} shows that the net flux $\mit\Gamma$
of the vector ${\bf g}$ through the whole magnetic surface is given by 
\begin{equation}
{\mit\Gamma}
=\int\mbox{div}\,{\bf g}\,dV
=\av{{{\bf g}}\.\na V}
=\prti{V}{a}\av{{{\bf g}}\.\na a}
\end{equation}
and the flux density
\begin{equation}
\gamma
=\frac{{\mit\Gamma}}{S_a}
=\frac{\av{{{\bf g}}\.\na a}}{\av{|\na a|}}
.\end{equation}

Assume now that a function of magnetic surface $F(a)$ describes a
quantity which satisfies the general diffusion equation 
\begin{equation}\label{dif1}
\prti{F}{t}+\av{\mbox{div}\,{\bf g}}
=\prti{F}{t}+\prti{\mit\Gamma}{V}
=S(a)
.\end{equation}
%which can also be represented in any of the forms
%\begin{equation}\label{locdif}
%\prti{F}{t}+\prti{\mit\Gamma}{V}=S(a)
%,\qquad\qquad
%\prti{F}{t}+\prti{a}{V}\prti{}{a}\lb(\prti{V}{a} \av{|\na a|}\gamma
%\rb)=S(a) 
%.\end{equation}
The local flux ${\bf g}$ can be taken in the form
\begin{equation}\label{Gloc}
{\bf g}(a,\theta)
=F(a)\widetilde{\bf v}(a,\theta)-\widetilde D(a,\theta)\na F(a)
.\end{equation}
On substitution in \req{dif1} one has
\begin{equation}\label{avdif}
\prti{F}{t}
=\prti{}{V}\lb( \av{(\na V)^2 \widetilde D}\prti{F}{V} - F\av{\na
V\.\widetilde{\bf v}} \rb) 
+S(a)
.\end{equation}
This result suggests the introduction of two functions of a single
argument $a$ 
\begin{equation}\label{Dv}
D(a)
=\frac{\av{(\na V)^2 \widetilde D}}{\av{(\na V)^2}}
=\frac{\av{(\na a)^2 \widetilde D}}{\av{(\na a)^2}}
,\qquad\quad
v(a)
=\frac{\av{\na V\.\widetilde{\bf v}}}{\av{|\na V|}}
=\frac{\av{\na a\.\widetilde{\bf v}}}{\av{|\na a|}}
.\end{equation}
These definitions, obviously, do not depend on the particular choice of
the magnetic surface label $a.$ 
Equation (\ref{avdif}) then reads
\begin{equation}\label{dif}
\prti{F}{t}
=\prti{a}{V}\prti{}{a}\lb[\prti{V}{a}\av{(\na a)^2}
\lb(D\prti{F}{a} - \frac{\av{|\na a|}}{\av{(\na a)^2}} v F\rb)
\rb]+S(a)
.\end{equation}
This presentation of the diffusion equation is employed in the Astra code.
The total flux and the average flux density on a magnetic surface are
\begin{equation}\label{Gamma}
{\mit\Gamma}(a)
=\prti{V}{a}
\lb(\av{|\na a|} v F - \av{(\na a)^2}D\prti{F}{a} \rb)
,\qquad\qquad
\gamma(a)
=v F - \frac{\av{(\na a)^2}}{\av{|\na a|}}D\prti{F}{a} 
.\end{equation}
If it is needed to introduce an ``effective'' diffusion coefficient as
implied by experimental observations then none of \req{Gamma} is
appropriate. 
The coefficient can be derived from \req{dif} as
\begin{equation}\label{Deff}
D_{\rm eff}(a)
=\frac{\dstyle\int\prti{F}{t}\,dV-\int S\,dV}
{\av{(\na V)^2}\prti{F}{V}}
=\frac{\dstyle\int\lb(\prti{F}{t}-S\rb)\prti{V}{a}\,da}
{\av{(\na a)^2}\prti{V}{a}\,\prti{F}{a}}
\end{equation}
which under the additional assumption $v=0$ coincides with $D(a)$ as
defined by \req{Dv}. 

The definitions (\ref{Dv}) are not unique. 
For instance, it is possible to write the average transport equation in
the form 
\begin{equation}\label{dif11}
\prti{F}{t}
=\prti{a}{V}\prti{}{a}\lb[\prti{V}{a}\av{|\na a|}
\lb(\widehat D\prti{F}{a} - v F\rb)
\rb]+S(a)
\end{equation}
rather than \req{dif} with the argument that $\prti{V}{a}\av{|\na a|}$ is 
the surface area and that the equation (\ref{dif11}) 
and the flux density $\gamma=vF-\widehat D\;\prt{F}/\prt{a}$ appear in
more ``cylindrical'' form. 
The resulting expression for the diffusion coefficient
$
\widehat D_{\rm eff}(a)
=\lb.\dstyle\int\lb(\prt{F}/\prt{t}-S\rb)dV\rb/\lb(S_a\;
\prt{F}/\prt{a}\rb)  
$
looks simpler and, probably, more natural, but it is dependent on 
a choice of the magnetic surface label ``{\it a}''. 
In the Astra code, the invariant definition (\ref{Deff}) is used.

We also introduce a local velocity, ${\bf u}_a$, of the $a(r,z)={\rm Const}$ surface as defined by the equation
\begin{equation}
\label{u}
\left.\prti{a}{t}\right|_{\vec r} 
+{\bf u}_a\cdot\na a = 0
.\end{equation}
The local velocity ${\bf u}_f$ of another flux surface labeled by $f$
(where $f=f(a,t)$) may be different from ${\bf u}_a$ but both are
obviously related as 
\begin{equation}\label{urel}
\left.\prti{f}{a}\right|_t
\left({\bf u}_a-{\bf u}_f\right)\cdot\na a
= \left.\prti{f}{t}\right|_a 
.\end{equation}
Hence $\left.\prt{f}/\prt{t}\right|_a$ expresses the relative motion of
the flux surface $f$ with respect to the flux surface $a.$ 
In other words, a time derivative of a surface function depends on the
assumption which quantity is taken as the `radial' coordinate and
considered to be fixed while computing $\prt/\prt t.$ 

\subsection{Formulae for plasma current and magnetic field}

Taking advantage of the toroidal symmetry and making use of the Maxwell
equations we can express two three-dimensional vectors, the magnetic
field {\bf B} and the current density {\bf j}, in terms of two scalar
functions $I(a)$ and $\Psi(a)$ (SI units are used): 
\begin{equation}\label{dfB}
\dstyle
{\bf B}
=I\nb\zeta+\frac{1}{2\pi}\left[\nb\Psi\times\nb\zeta\right]
,\end{equation}
\begin{equation}\label{dfj}
\dstyle
{\bf j}
=-\frac{\nb\zeta}{2\pi\mu_0}r^2{\rm div}\frac{\nb\Psi}{r^2}
+\frac{1}{\mu_0}\left[\nb I\times\nb\zeta\right]
.\end{equation}
Taking the scalar product of equations (\ref{dfB}) and (\ref{dfj}) with
$\nb\theta$ and integrating over the volume within magnetic surface
using \req{dfV}, we obtain\footnote{In what follows we use exclusively
the quantity $\psi=-\Psi=-2\pi rA_\zeta$ with $A_\zeta$ being the
toroidal component of the magnetic vector potential.}: 
\begin{equation}
\label{dfpsi}
\psi
=-\Psi
=\frac{1}{2\pi}\int\limits_{\mbox{\msym V}}{\bf B}\cdot\nb\theta\, d^3x
\equiv
\int\limits_{{\mbox{\msym S}}_\theta}{\bf B}\cdot d{\bf S}_\theta
,\end{equation}
\begin{equation}
\label{dfI}
I
=R_0 B_0-\frac{\mu_0}{4\pi^2}\int\limits_{\mbox{\msym V}}{\bf
j}\cdot\nb\theta\, d^3x 
\equiv
R_0 B_0-\frac{\mu_0}{2\pi}
\int\limits_{{\mbox{\msym S}}_\theta}{\bf j}\cdot d{\bf S}_\theta
.\end{equation}
Here $R_0$ is the distance from the axis of the torus to a fixed point 
inside the plasma.
Usually, one chooses $R_0$ as the geometric center of the vacuum chamber.
$B_0$ is the vacuum magnetic field at $r=R_0$. 
While $R_0$ is constant, the magnetic field at this point may be a
function of the time, $B_0=B_0(t)$. 
The additive constants in \eqs\ref{dfpsi}),~(\ref{dfI}) are chosen so
that the function $\psi=0$ at the magnetic axis and is physically
the poloidal magnetic flux, while the function $I$ is related to the
poloidal current inside a magnetic surface. 
As long as no current flows outside the plasma, it follows that $I\equiv
R_0B_0$ everywhere in space between the plasma edge and the toroidal
field coils.  
Equations \rf{dfpsi} and \rf{dfI} can be viewed as definitions for
functions $\psi$ and $I$, respectively. We introduce also the
dimensionless quantity 
\begin{equation}
\label{dfJ}
J\df\frac{I}{R_0 B_0}
.\end{equation}
As long as poloidal (diamagnetic) plasma current is much smaller than
the current in toroidal field coils, the quantity $J$ is very close to
unity inside the plasma and $J=1$ outside it. 

Both functions defined in \eqs\ref{dfpsi}), (\ref{dfI}) are so called
surface functions that means that they depend on space
coordinates only through the variable $a$. 
In evolving plasmas, all surface functions may depend also on
time $t,$ e.g., $V=V(a,t)$, $\psi=\psi(a,t)$. 
Each of them can be used as radial coordinates instead of $a$. 
It is convenient to introduce two further surface functions: 
\begin{equation}
\label{dfPhi}
\Phi
\df\int\limits_{{\mbox{\msym S}}_\zeta}{\bf B}\cdot d{\bf S}_\zeta
=\frac{1}{2\pi}\int\limits_{\mbox{\msym V}}\left({\bf
B}\cdot\nb\zeta\right) d^3x 
=\frac{1}{2\pi}\int\limits_{\mbox{\msym V}}\frac{I}{r^2} d^3x
,\qquad
\rho\df\sqrt{\Phi\over\pi B_0}
.\end{equation}
The first of them is the toroidal magnetic flux $\Phi.$ The second
$\rho,$ has the dimensionality of length and represents an effective
minor radius. 

As seen from \req{dfPhi} the quantities $I$ and $\Phi$ can be expressed
in terms of each other. 
On differentiation of \req{dfPhi} we have
\begin{equation}
\label{ident1}
\prti{\Phi}{V}
=\frac{I}{2\pi}\av{\frac{1}{r^2}}
\qquad\mbox{or}\qquad
4\pi^2\rho R_0
=J\prti{V}{\rho}\av{\frac{R_0^2}{r^2}}
\end{equation}
which in large aspect ratio (where $\av{R_0^2/r^2}\approx 1$) and at low
plasma pressure (where $J\approx 1$) reduces to
$\prt{V}/\prt{\rho}\approx 4\pi^2\rho R_0$ or 
$V\approx 2\pi^2\rho^2R_0$ showing that the variable $\rho$ can indeed
be viewed as an effective (cylinder-like) minor radius of a magnetic
surface. 
In what follows, we use $\rho$ as the main radial variable and assume
that all surface functions are functions of $\rho$ and $t$,
i.e. $\psi=\psi(\rho,t)$, $V=V(\rho,t)$ and so on. 
It means that time derivatives are understood as being computed with
$\rho$ kept constant: 
\begin{equation}\label{doverdt}
\prti{f}{t}
\df\prti{f(\rho,t)}{t}
\equiv
\left.\prti{f}{t}\right|_\rho
.\end{equation}
unless declared otherwise. Time derivatives taken at
constant $\Phi$ also play a significant role and are related to the
standard time derivative of \req{doverdt} by the following relation:  
\begin{equation}
\label{dtphi}
\left.\prti{f}{t}\right|_\Phi
=\left.\prti{f(\rho(\Phi,t),t)}{t}\right|_\Phi
=\left.\prti{f}{t}\right|_\rho
+\prti{f}{\rho}\left.\prti{\rho}{t}\right|_\Phi
=\prti{f}{t}
-\frac{\rho\dot B_0}{2B_0}\prti{f}{\rho}
.\end{equation}
Some flux  surface averages are also included in the transport equations 
and we introduce the following special notations: 
\begin{equation}\label{metr}
V'=\prti{V}{\rho}
,\quad
G_1\df\av{\left(\nb \rho\right)^2}
,\quad
G_2\df\frac{V'}{4\pi^2}\av{\left(\frac{\nb \rho}{r}\right)^2}
,\quad
G_3\df\av{\frac{R_0^2}{r^2}}
\equiv\frac{4\pi^2\rho R_0}{JV'}
\end{equation}
where all $G_{\rm i}$ are dimensionless.

For some applications the poloidal and toroidal components of the
magnetic field {\bf B} are required. They can be found from \req{dfB} 
\begin{equation}
\label{Bpol}
B_{\rm pol}
=\frac{\left|\nb\rho\right|}{2\pi r}\prti{\psi}{\rho}
,\qquad\qquad
B_{\rm tor}=\frac{I}{r}
.\end{equation}
Both components are poloidally dependent. 
We introduce also the average cylinder-like poloidal field $B_{\rm p}$,
the rotational transform $\mu$ and the safety factor $q$ as 
\begin{equation}\label{dfmu}
B_{\rm p}
\df\frac{1}{2\pi R_0}\prti{\psi}{\rho}
,\qquad\qquad
%Definition for \mu
\iota\df\mu\df
\prti{\psi}{\Phi}
=\frac{1}{2\pi B_0\rho}\prti{\psi}{\rho}
=\frac{B_{\rm p} R_0}{B_0\rho}
,\qquad\qquad
q\df
\frac{1}{\mu}
=\prti{\Phi}{\psi}
.\end{equation}
Averaging \req{dfj} and using \req{avdiv} we find the toroidal current
density 
\begin{equation}
\label{dfjtor}
j_{\rm tor}
\df R_0\av{\bf j\cdot\nb\zeta}
=\frac{2\pi R_0}{\mu_0 V'}\prti{}{\rho}
\left(G_2\prti{\psi}{\rho}\right)
=\frac{J G_3}{2\pi\mu_0\rho}\prti{}{\rho}
\left(G_2\prti{\psi}{\rho}\right)
\end{equation}
and the parallel current density
\begin{equation}\label{dfjp}\dstyle
%Definition for the parallel current
j_\pll
\df \frac{\av{\bf j\cdot\bf B}}{B_0}
\dstyle
=\frac{2\pi R_0}{\mu_0 V'}J^2
\prti{}{\rho}
\left(G_2 J^{-1}\prti{\psi}{\rho}\right)
.\end{equation}
Integration of \req{dfjtor} gives different relations for the toroidal
current 
\begin{equation}
\label{Ipl}
\begin{array}{rl}\dstyle
I_{\rm pl}
\df\int\limits_{\mbox{\msym S}_\zeta}{\bf j}\cdot d{\bf S}_\zeta
&\dstyle
=\frac{1}{2\pi}
\int\limits_{\mbox{\msym V}}\left({\bf j}\cdot\nabla\zeta\right) d^3x
=\frac{1}{2\pi R_0}\int\limits_0^\rho V' j_{\rm tor} d\rho
=\frac{J}{2\pi R_0}\int\limits_0^\rho \frac{V'}{J^2} j_\| d\rho
\bigskip \\ &\dstyle
=\frac{G_2}{\mu_0}\prti{\psi}{\rho}
=\frac{2\pi B_0}{\mu_0}\rho G_2\mu
=\frac{2\pi R_0}{\mu_0}G_2 B_{\rm p}
\end{array}
\end{equation}
and
\begin{equation}\label{jlon}
j_{\rm tor}
=2\pi R_0\prti{I_{\rm pl}}{V}
,\qquad
j_\|
=2\pi R_0J^2\prti{}{V}\lb(\frac{I_{\rm pl}}{J}\rb)
=J^2\prti{}{V}\int\limits_0^V\frac{j_{\rm tor}}{J}dV
.\end{equation}
Also the formula for the poloidal field energy, which follows from 
\eqs\ref{metr})-(\ref{Bpol}), is useful 
\begin{equation}
\label{WI}
W_{\rm I}
=\int\limits_{\mbox{\msym V}}\frac{{\bf B}_{\rm pol}^2}{2\mu_0}d^3x
=\frac{1}{2\mu_0}\int\limits_0^\rho 
\left(\prti{\psi}{\rho}\right)^2 G_2 d\rho
.\end{equation}
The internal inductance $L_{\rm i}^W$ of the plasma current can be found
from \eqs \ref{Ipl}), (\ref{WI}) as 
$L_{\rm i}^W=2W_{\rm I}/I_{\rm pl}^2$. 
The frequently used dimensionless quantity which is related to the
internal inductance per unit length, $l_{\rm i}$, is given by $l_{\rm
i}=\dstyle\frac{4\pi}{\mu_0}\frac{L_{\rm i}^W}{2\pi R_0}=\frac{2L_{\rm
i}^W}{\mu_0 R_0}$. 
Sometimes the internal inductance is defined as ${L_{\rm
i}^\psi=[\psi(a_{\rm pl})-\psi(0)]/I_{\rm pl}}$ which is obviously
different from $L_{\rm i}^W$. 

\subsection{Equilibrium equation}

Plasma configuration in a tokamak is determined by the Grad-Shafranov
equilibrium equation 
\begin{equation}
%G-Sh equation 
\label{GSh}
\Delta^*\psi
=r^2{\rm div}\frac{\nb\psi}{r^2}
%=2\pi\mu_0r^2\left({\bf j}\cdot\nabla\zeta\right)
=-4\pi^2\left(\mu_0 r^2\prti{p}{\psi}+I\prti{I}{\psi}\right)
%=-\frac{2\pi}{R_0 B_{\rm p}}\left(\mu_0 r^2\prti{p}{\rho}
%+I\prti{I}{\rho}\right)
.\end{equation}
The first term on the \rhs  contains the plasma pressure $p=p(\rho,t)$
where the contributions from all plasma species are included. The second
term describes the poloidal diamagnetic current, which can be readily
expressed in terms of the average toroidal current density. 
Eliminating $\Delta^*\psi$ from \eqs\ref{dfB}), (\ref{dfj}) and
(\ref{GSh}) we find
\begin{equation}\label{j1}
{\bf j}
=-2\pi r^2\prti{p}{\psi}\na\zeta-\frac{2\pi}{\mu_0}\prti{I}{\psi}{\bf B}
.\end{equation}
As mentioned above, the transport equations include parallel component
of the current density. 
So we need to represent the quantity $\prt I/\prt\psi$ in terms of the
parallel current density $j_\pll$. 
A scalar product of \eqs\ref{dfB}) and (\ref{j1}) gives
\begin{equation}\label{j2}
j_\pll
=-\frac{2\pi}{B_0}\lb(I\prti{p}{\psi}
+\frac{\av{B^2}}{\mu_0}\prti{I}{\psi}\rb)
,\end{equation}
where
\begin{equation}
\frac{\av{B^2}}{B_0^2}
=G_3 J^2+\frac{4\pi^2 G_2\rho^2 \mu^2}{V'}
.\end{equation}
We can now express the \rhs of \req{GSh} in terms of functions which are
provided by the transport equations 
\begin{equation}\label{GSp}
\begin{array}{rl}\dstyle
\Delta^*\psi
&\dstyle
%=\frac{2\pi\mu_0 R_0 J}{\av{B^2/B_0^2}}j_\pll
%+4\pi^2\mu_0R_0^2\lb(\frac{J^2}{\av{B^2/B_0^2}}
%-\frac{r^2}{R_0^2}\rb)\prti{p}{\psi}
=\frac{2\pi\mu_0 R_0 J}{\av{B^2/B_0^2}}j_\pll
+\frac{2\pi\mu_0R_0^2}{B_0\rho\mu}\lb(\frac{J^2}{\av{B^2/B_0^2}}
-\frac{r^2}{R_0^2}\rb)\prti{p}{\rho}
\bigskip \\ &\dstyle
=2\pi\mu_0 R_0\lb[
\frac{J}{\av{B^2/B_0^2}}\lb(j_\pll
+\frac{R_0 J}{B_0\rho\mu}\prti{p}{\rho}\rb)
-\frac{r^2}{B_0 R_0\rho\mu}\prti{p}{\rho}\rb]
.\end{array}
\end{equation}

The equilibrium equation (\ref{GSp}) depends on $t$ parametrically but
not explicitly.  
Physically this means that we assume that the plasma is always in 
equilibrium, and we do not consider relaxation processes. 
The justification for this is that the relaxation to equilibrium is
several orders of magnitude faster than all transport processes. 
We consider the coupling of this equation to the time-dependent transport
equations in Section 3.11. 

\subsection{Toroidal electric field}

For the toroidal vortex electric field ${\bf E}$ we have local
(poloidally dependent) relation 
\begin{equation}
\label{vecE}
{\bf E}
=-\frac{1}{2\pi}\left.\prti{\Psi}{t}\right|_{\bf r}\nabla\zeta
=\frac{1}{2\pi}\left.\prti{\psi}{t}\right|_{\bf r}\nabla\zeta
.\end{equation}
The toroidal loop voltage 
$U_{\rm tor}$
is usually understood as the quantity measured by a fixed toroidally
symmetric loop of a constant major radius $\br_l=(r_l,z_l)$ which is
given by 
\begin{equation}
\label{dfUt}
U_{\rm tor}=2\pi r\left({\bf E}\cdot r\nabla\zeta\right)
=\dstyle\left.\prti{\psi}{t}\right|_{\br_l}
=-\lb.\lb({\bf u}_\psi\cdot\nabla\psi\rb)\rb|_{\br_l}
\end{equation}
where ${\bf u}_\psi$ a local velocity of the constant-$\psi$ surface
(see \req{u}). Different kinds of motion contribute to ${\bf u}_\psi$
and, hence, to $U_{\rm tor}$.  

First of all, we have 
\begin{equation}
\label{UPhi}
U_{\rm tor}
=\left.\prti{\psi(\Phi,t)}{t}\right|_{\br_l}
=\left.\prti{\psi}{t}\right|_\Phi
+\mu\left.\prti{\Phi}{t}\right|_{\br_l}
.\end{equation}
It can be shown \cite{HH} that the first term on the \rhs  of \req{UPhi}
is related to an average longitudinal electric field $E_\pll$ on the
magnetic surface $\Phi$ 
\begin{equation}
\label{dfEp}
U_\pll
\df
2\pi R_0 E_\pll
\df
\frac{2\pi R_0}{B_0}\av{\bf E\cdot B}
=\left.JG_3\prti{\psi}{t}\right|_\Phi
=JG_3\lb(\left.\prti{\psi}{t}\right|_\rho-\pi\rho^2\mu\dot B_0\rb)
.\end{equation}
In ideally conducting plasmas $E_\pll=0$ and 
%$\left.\prt\psi/\prt t\right|_\Phi=0$ 
\req{dfEp} describes freezing of the fluxes $\psi$ and $\Phi$ in one
another. 
Using Eqs.~(\ref{urel}), (\ref{ident1}) and (\ref{dfUt}) we rewrite
\req{UPhi} as 
\begin{equation}
\label{psivsPhi}
\mu\left({\bf u}_\psi
-{\bf u}_\Phi\right)\cdot\nb\Phi
+\frac{V'}{4\pi^2 R_0\rho}U_\pll
=0
.\end{equation}
We conclude that the poloidal flux $\psi$ moves through the toroidal
flux $\Phi$ only in proportion to its resistive dissipation within the
flux surface. 
The quantity $U_\pll$ (or $E_\pll$) is a measure of such dissipation. 
It is the quantity $U_\pll$ which directly participates in and is
provided by a transport code. 
It describes irreversible diffusive flow of the poloidal flux $\psi$
through the toroidal flux $\Phi.$ 

Making use of \req{dfEp} we introduce another flux surface quantity,
$U_{\rm pl},$ which is the voltage as seen in the coordinate system
moving together with a $\rho$--surface: 
\begin{equation}
\label{Upl}
U_{\rm pl}
=\left.\prti{\psi}{t}\right|_\rho
=\frac{U_\pll}{JG_3}
+\pi\rho^2 \mu\dot B_0
%=U_{\rm tor}+2\pi \rho\mu B_0\left({\bf u}_\rho\cdot\nb\rho\right)
%=U_{\rm tor}-2\pi \rho\mu B_0\left.\prti{\rho}{t}\right|_{\br_l}
=\frac{U_\pll}{JG_3}
+2\pi\rho\mu B_0\left({\bf u}_\rho-{\bf u}_\Phi\right)\cdot\nb\rho
,\end{equation}
where \req{urel} and the definition \req{dfPhi} have also been used. 
The last equality in \req{Upl} shows that the relative motion of the
$\Phi$--surface with respect to the $\rho$--surface appears only due to
variation of the external toroidal magnetic field $B_0$. 
In most present day tokamaks the vacuum magnetic field can be viewed as
time-independent, $\dot B_0=0,$ and both flux surface labels,
$\rho$ and $\Phi,$ are physically equivalent. 
Usually, the coefficient in front of $U_\pll$ in \req{Upl} is close to
unity and, when $\dot B_0=0$, $U_{\rm pl}$ is close to $U_\pll$. 
However, the difference between $U_\pll$ and $U_{\rm pl}$ can become
significant in a small aspect ratio tokamak and at high plasma pressure. 
For practical applications, the quantity $U_{\rm pl}$ is more convenient
than $U_\pll$ because, {\em in steady state}, 
$U_{\rm pl}$ does not vary over the minor radius $\rho:$ 
\begin{equation}
\label{Uplss}
\prti{U_{\rm pl}}{\rho}=0
.\end{equation}

We consider now the second term on the \rhs  of \req{UPhi}. 
First of all, we note that the pressure of the toroidal magnetic field
in a tokamak is much higher than the plasma pressure. 
For this reason the toroidal flux $\Phi$ cannot be changed noticeably
by any internal processes in the plasma. 
Consequently, the second term on the \rhs  of \req{UPhi} can
only become appreciable when the external magnetic fields vary. 
For instance, varying of the poloidal control field causes
plasma movement as a whole (compression in major radius), while increase
of the toroidal field results in compression in minor radius.
When any of these movements is made faster than characteristic transport
times then the plasma is compressed adiabatically.
We can further consider the equality 
\begin{equation}
\label{drho}
-{\bf u}_\rho\cdot\nb\rho
=\left.\prti{\rho}{t}\right|_{\br_l}
=\prti{\rho}{V}\left.\prti{V}{t}\right|_{\br_l}
+\left.\prti{\rho}{t}\right|_V
=\prti{\rho}{V}\lb(\left.\prti{V}{t}\right|_{\br_l}
-\left.\prti{V}{t}\right|_\rho\rb)
\end{equation}
relating the ``magnetic'' variable $\rho$ to the ``geometric'' variable
$V$. We conclude that ${\bf u}_\rho$ comprises the plasma motion in
space as a whole and a motion of the $\rho$--surface because of changes
in the equilibrium (first and second
terms in the brackets on the \rhs of \req{drho}, respectively).  

In general case, using \eqs\ref{UPhi})--(\ref{drho}) we can write the
toroidal loop voltage as 
\begin{equation}\label{Utor1}
U_{\rm tor}
=\frac{1}{JG_3}U_\pll
+2\pi \rho\mu B_0\lb.\prti{\rho}{t}\rb|_V
+\pi\rho^2 \mu\dot B_0
+\frac{2\pi \rho\mu B_0}{V'}\left.\prti{V}{t}\right|_{\br_l}
.\end{equation}
All terms here have clear physical significance. 
The \lhs  of \req{Utor1} is the loop voltage as measured in an experiment.
The first term on the \rhs gives a resistive component. 
The second term describes evolution of the plasma surface due to a
variation in plasma para- or dia-magnetism. 
It is of order $d\be_{\rm tor}/dt$ and in a tokamak is usually
negligibly small. 
The remaining two terms could be associated with the adiabatic compression
in minor and major radii, respectively. 

In order to express the measured loop voltage $U_{\rm tor}$ in terms of
the calculated quantity $U_{\rm pl}$ we rewrite \req{Utor1} as
\begin{equation}\label{Utor}
U_{\rm tor}
=U_{\rm pl}
+2\pi \rho\mu B_0\left.\prti{\rho}{t}\right|_V
+\frac{2\pi \rho\mu B_0}{V'}\left.\prti{V}{t}\right|_{\br_l}
\end{equation}
and evaluate it at some point $\br_l=\br_B$ assuming that the point
$\br_B$ belongs to the edge flux surface with a fixed (for instance, 
due to the boundary conditions) volume $V=V_B$ so that 
$\prt V_B/\prt t=0.$ 
\begin{equation}\label{Uplb}
U_{\rm plB}
=U_{\rm tor}\Bigr|_{\rho_B}
=U_{\rm pl}(\rho_B)
+2\pi B_0\rho_B\mu(\rho_B)\prti{\rho_B}{t}
=\lb.\prti{\psi}{t}\rb|_{\rho_B}
+2\pi B_0\rho_B\mu(\rho_B)\prti{\rho_B}{t}
%+\frac{2\pi \rho\mu B_0}{V'}\left.\prti{V_B}{t}\right|_{\br_B}
.\end{equation}
The last term on the \rhs of \req{Uplb} is nonzero (although quite small)
because the toroidal flux inside the plasma changes with any change in
the profiles of plasma pressure or current density which gives rise to
$\rho_B=\rho_B(t).$ 
This term is provided by a solution of the Grad-Shafranov equation.
Evaluation of $U_{\rm tor}$ at arbitrary point $\br_l=(r_l,z_l)$ outside
the plasma requires solution of the equilibrium equation in the gap
between the plasma boundary and the position of measurement $\br_l.$ 
This problem is not considered here and, in what follows, we assume that
the edge loop voltage is calculated according to \req{Uplb}.

\subsection{Joule heating and Ohm's law}

Dissipation of the magnetic field energy within a flux surface
is defined by 
\begin{equation}\label{QJ1}\hspace*{-10pt}
\begin{array}{rl}\ds
\int\limits_{\mbox{\msym V}} \left({\bf j\cdot E}\right) d^3x
&\!\!\!\dstyle
=\frac{1}{4\pi^2\mu_0}\int\limits_{\mbox{\msym V}} 
\left.\prti{\psi}{t}\right|_\br
{\rm div}\frac{\nabla\psi}{r^2}\;d^3x
%\bigskip\\ &\dstyle
=\frac{1}{4\pi^2\mu_0}\int\limits_{\mbox{\msym V}} 
\lb(\left.\prti{\psi}{t}\right|_\rho
+\prti{\psi}{\rho}\left.\prti{\rho}{t}\right|_\br\rb)
{\rm div}\frac{\nabla\psi}{r^2}\;d^3x
%\bigskip\\ &\!\!\!\dstyle
%=\frac{1}{4\pi^2\mu_0}\int\limits_{\mbox{\msym V}} 
%\lb[U_{\rm pl}-2\pi B_0\rho\mu\left({\bf u}_\rho\cdot\nb\rho\right)\rb]
%{\rm div}\frac{\nabla\psi}{r^2}\;d^3x
\bigskip\\ &\!\!\!\dstyle
=\int\limits_0^\rho U_{\rm pl}\prti{I_{\rm pl}}{\rho}d\rho
-\frac{B_0}{2\pi\mu_0}\int\limits_{\mbox{\msym V}} 
\rho\mu\left({\bf u}_\rho\cdot\nb\rho\right)
{\rm div}\frac{\nabla\psi}{r^2}\;d^3x
.\end{array}
\end{equation}
The first term on the \rhs of \req{QJ1} can be transformed as
\begin{equation}\label{QJ2}
%\begin{array}{rl}
Q_{\rm J}
%&\dstyle
%=\frac{1}{4\pi^2\mu_0}\int\limits_{\mbox{\msym V}} 
%\left.\prti{\psi}{t}\right|_\rho
%{\rm div}\frac{\nabla\psi}{r^2}\;d^3x
=\int\limits_0^\rho U_{\rm pl}\prti{I_{\rm pl}}{\rho}d\rho
%\bigskip\\ &\dstyle
=I_{\rm pl}U_{\rm pl}
-\frac{1}{\mu_0}
\int\limits_0^\rho
\prti{\psi}{\rho}\prtii{\psi}{t}{\rho} G_2 d\rho
=I_{\rm pl}U_{\rm pl}
-\frac{1}{2\mu_0}
\prti{}{t}
\int\limits_0^\rho
\left(\prti{\psi}{\rho}\right)^2 G_2 d\rho
%.\end{array}
\end{equation}
which on account of the definition (\ref{WI}) can be represented as
the conservation law for the energy of the poloidal magnetic field in
the coordinate system moving together with a $\rho$--surface 
\begin{equation}\label{WIcons}
\prti{W_{\rm I}}{t}
=I_{\rm pl}U_{\rm pl}-Q_{\rm J}
.\end{equation}
Here $I_{\rm pl}U_{\rm pl}$ is the Pointing flux relative to the
$\rho={\rm Const}$ flux surface. 

We can also rewrite \req{QJ2} as
\begin{equation}
\label{QJ}
Q_{\rm J}
=\int\limits_0^\rho U_{\rm pl}\prti{I_{\rm pl}}{\rho}d\rho
=\frac{1}{2\pi R_0}\int\limits_0^\rho U_{\rm pl}j_{\rm tor}V'd\rho
=\frac{1}{2\pi R_0}\int\limits_0^V U_{\rm pl}j_{\rm tor}dV
,\end{equation}
thus finding the density of magnetic energy dissipation as
\begin{equation}\label{PJ}
P_{\rm J}=\frac{1}{2\pi R_0}U_{\rm pl}j_{\rm tor}
.\end{equation}
This expression does not take into account the work done by the moving
surface against the plasma pressure gradient described by the second
term on the \rhs of \req{QJ1}. 
As discussed in the previous section this contribution is usually small
unless the fast processes as the adiabatic compression are involved
~\cite{HH}. 

The longitudinal Ohm law is assumed to have the form
\begin{equation}\label{Ohm}
j_\pll
=\sigma_\pll E_\pll+j_{\rm BS}+j_{\rm CD}
\end{equation}
which together with relations \req{dfjp} and \req{dfEp} gives a transport
equation for the poloidal flux $\psi$ considered in the next section.
The averages of the bootstrap current density, ${\bf j}_{\rm BS}$, and
the density of the current driven by external sources, ${\bf j}_{\rm
CD}$, are determined similar to $j_\pll$ in \req{dfjp} 
\begin{equation}\label{jCD}
j_{\rm BS}=\dstyle\frac{1}{B_0}\av{{\bf j}_{\rm BS}\cdot{\bf B}}
,\qquad\qquad
j_{\rm CD}=\dstyle\frac{1}{B_0}\av{{\bf j}_{\rm CD}\cdot{\bf B}}
.\end{equation}
%As can be shown, only a part 
%\begin{equation}\label{POH}
%P_{\rm OH}=\sigma_\pll E_\pll^2
%\end{equation}
%of the total power $P_{\rm J}$ lost by the poloidal field goes
%irreversibly to the bulk electron heating. There is another part, which
%goes directly to non-Maxwellian particles and is spent for their
%accelerating or decelerating. Under some conditions, this power can be
%returned to the poloidal magnetic field system. 


%\vfill\eject

%\section{Tokamak simulation model}

\subsection{Transport equations}

The tokamak simulation code Astra like other codes comprises of a
system of 1D diffusion equations for densities and temperatures of
different plasma components, a 2D equilibrium equation, and, 
a variety of other modules to describe additional heating, current
drive and other non-diffusive processes in tokamak plasmas. A set of the
transport equations for a tokamak is derived in \cite{HH}. In this
section, we discuss the set of equations which are  implemented in the
Astra code which differs from that presented in \cite{HH} by allowance
for varying toroidal field $B_0$ and, consequently, adiabatic
compression in minor radius.  

The basic set of transport equations in the Astra code includes
equations for the electron density $n_e,$ electron temperature, $T_e,$
ion temperature, $T_i,$ and the poloidal flux~$\psi$ 
\begin{equation}\label{main}
\hspace*{-5mm}\lb\{
\begin{array}{l}\dstyle
\frac{1}{V'}\left(\prti{}{t}
-\frac{\dot B_0}{2B_0}\prti{}{\rho}\rho
\right)\left(V'n_e\right)
+\frac{1}{V'}\prti{}{\rho}
\Gamma_e
=S_e
,\bigskip\\ \dstyle
\frac{3}{2}\left(V'\right)^{-5/3}
\left(\prti{}{t}
-\frac{\dot B_0}{2B_0}\prti{}{\rho}\rho
\right)\left[
\left(V'\right)^{5/3}n_e T_e
\right]
+\frac{1}{V'}\prti{}{\rho}
\left(q_e+\frac{5}{2}T_e\Gamma_e
\right)
=P_e
,\bigskip\\ \dstyle
\frac{3}{2}\left(V'\right)^{-5/3}
\left(\prti{}{t}
-\frac{\dot B_0}{2B_0}\prti{}{\rho}\rho
\right)\left[
\left(V'\right)^{5/3}n_i T_i
\right]
+\frac{1}{V'}\prti{}{\rho}
\left(q_i+\frac{5}{2}T_i\Gamma_i
\right)
=P_i
%,\bigskip\\ \dstyle
%\frac{2\pi\rho}{V'}\sigma_\pll\left(\prti{\psi}{t}-\frac{\rho\dot B_0}{2B_0}\prti{\psi}{\rho}\right)=\frac{I^2}{2\pi\mu_0V'B_0}\prti{}{\rho}\left[\frac{V'}{I}\av{\left(\frac{\nb\rho}{r}\right)^2}\prti{\psi}{\rho}\right]-j_{BS}-j_{CD}
,\bigskip\\ \dstyle
%\label{Ohm}
\sigma_\pll
\left(\prti{\psi}{t}
-\frac{\rho\dot B_0}{2B_0}\prti{\psi}{\rho}
\right)
=\frac{J^2R_0}{\mu_0\rho}\prti{}{\rho}
\left(\frac{G_2}{J}
\prti{\psi}{\rho}
\right)
-\frac{V'}{2\pi\rho}\left(j_{BS}+j_{CD}\right)
.\end{array}\rb.
\end{equation}
If the main ion density $n_i$ and its flux $\Gamma_i$ are not explicitly
defined (see Section 4) then it is assumed that $n_i=n_e/Z_i$ and
$\Gamma_i=\Gamma_e/Z_i.$  
(Note that this approximation neglects impurities.) 
In line with the discussion in Section 2.1 all fluxes included in
\req{main}, i.e. the electron flux $\Gamma_e$, the electron heat flux
$q_e$ and the ion heat flux $q_i$, are considered as total fluxes
through a flux surface $\rho=Const.$ In order to close the system of
equations the fluxes should be expressed in terms of thermodynamic
forces taken as derivatives with respect to $\rho$. In a simulation, the
fluxes as well as the conductivity $\sigma_\pll$ and the bootstrap
current density $j_{BS}$ can either be taken from a transport theory or
from experimental observations in a tokamak. Independently of the origin
of the formulae used for the representation of the fluxes, it is assumed
in the Astra code that the transport matrix has the following form 
\begin{equation}\label{fluxes}
\left(
\begin{array}{c}\dstyle
\frac{\Gamma_e}{n_e}
\medskip\\ \dstyle
\frac{q_e}{n_eT_e}
\medskip\\ \dstyle
\frac{q_i}{n_iT_i}
\medskip\\ \dstyle
V'G_1\frac{\mu_0 j_{BS}}{B_p}
\end{array}
\right)
=-V'G_1
\left(
\begin{array}{cccc}
D_n & D_e & D_i & D_E
\bigskip\medskip\\ \dstyle
\chi_n^e & \chi_e & \chi_i^e & \chi_E^e
\bigskip\medskip\\ \dstyle
\chi_n^i & \chi_e^i & \chi_i & \chi_E^i
\bigskip\medskip\\ \dstyle
C_n & C_e & C_i & 0
\end{array}
\right)
\cdot
\left(
\begin{array}{c}\dstyle
\frac{1}{n_e}\prti{n_e}{\rho}
\medskip\\ \dstyle
\frac{1}{T_e}\prti{T_e}{\rho}
\medskip\\ \dstyle
\frac{1}{T_i}\prti{T_i}{\rho}
\medskip\\ \dstyle
\frac{E_\pll}{B_p}
\end{array}
\right)
.\end{equation}
All the coefficients in the upper left third-order minor of the
transport matrix have dimensions $[m^2/s]$ while
the right column and the bottom row are dimensionless. 
The average poloidal magnetic field $B_p$ is given by \req{dfmu}. 

In practice it is rare to use all the terms on the \rhs of the transport
system of equations \req{main} and the full transport matrix
\req{fluxes} are used in transport modeling.
Moreover, often the transport set of four equations \req{main} is
excessive. 
For instance, computation of electron density $n_e$ can often be
replaced with experimentally measured data. 
The Astra code provides an easy way to retain only those terms and
equations of \eqs\ref{main}),~(\ref{fluxes}) which are necessary for the
specific problem under consideration. 
This is coded as a set of special instructions and will be discussed
further in Section~5.2. 
On the other hand, the four transport equations of \req{main} can be
insufficient for some problems of interest, such as minority ion heating
or helium ash removal. In such cases, the code can be supplemented with
additional diffusion equations as will be described in Section~3.12. 

In addition to the main set of the transport equations (\req{main}) the
Astra code provides a variety of other modules for the description of
different processes in tokamak plasmas and to compute the transport
coefficients in \eqs{\ref{main}})-(\ref{fluxes}).
Processes modeled by these modules include additional heating and
current drive, minority species behavior, the plasma periphery and SOL
plasma, MHD stability analysis etc.

\subsection{Sources and sinks}

The source of electrons $S_e$ in the first of \eqs\ref{main}) is 
\begin{equation}\label{S}
S_e=s_{ion}^{(e)}Nn_e+s_{ion}^{(i)}Nn_i-s_{rec}n_in_e
\end{equation}
where $N$ is the density of neutral atoms of the working gas, $s_{rec}$
is the recombination rate, and $s_{ion}^{(e)}$ and $s_{ion}^{(i)}$ are
the rates of the impact ionization by electrons and ions, respectively. 
The sources of the electron $P_e$ and of the ion $P_i$ energy are
defined as
\begin{equation}
\begin{array}{l}\label{P}\dstyle
P_e
=P_{OH}-P_{\Gamma}-P_{ei}-P_e^{RAD}-P_e^N+P_e^H+P^{FUS}
,\bigskip\\ \dstyle
P_i=P_{\Gamma}+P_{ei}+P_i^N+P_i^H
,\end{array}
\end{equation}
include the Ohmic heating $P_{OH}$ which was defined in \req{PJ}, the
power lost as radiation $P_e^{RAD}$, the powers of auxiliary heating
of electrons $P_e^H$ and ions $P_i^H$, heat exchange between electron
and ion components 
\begin{equation}\label{Pei}\dstyle
P_{\Gamma}
=\av{(\nb\rho)^2}\frac{\Gamma_e}{n_e}\prti{\left(n_iT_i\right)}{\rho}
,\qquad
P_{ei}=\frac{3m_e}{m_i}\frac{n_e}{\tau_e}(T_e-T_i)
\end{equation}
and losses due to the atomic processes
\begin{equation}
\begin{array}{l}\label{PN}\dstyle
P_e^N
=\left(s_{ion}^{(e)}N{\cal E}_{ion}+s_{rad}N{\cal E}_1
+\tstyle\frac{3}{2}s_{rec}n_iT_e\right)n_e
,\bigskip\\ \dstyle
P_i^N
=\tstyle\frac{3}{2}\left(s_{ion}NT_N+s_{cx}N(T_N-T_i)
-s_{rec}n_eT_i\right)n_i
,\end{array}
\end{equation}
where $T_N$ is the temperature of the neutral atoms, $s_{cx}$ and
$s_{rad}$ are  the charge exchange and radiation rates, respectively,
${\cal E}_{ion}=13.6$ eV, ${\cal E}_1=10.2$ eV. 

Although the bootstrap current $j_{BS}$ is expressed by \req{fluxes} in
terms of thermodynamic forces, the bootstrap current together with the
driven current $j_{CD}$ can also be viewed as sources of poloidal flux. 
The driven current density $j_{CD}$ as well as $N$, $T_N$, $P_e^H$ and
$P_i^H$ must be provided by separate modules which are included in the
code. 
Sawtooth oscillations and plasma behavior in a divertor region are
further examples of processes which require separate treatment. 
Modules for these processes are also included in the code, but
detailed discussion of all these modules falls outside the scope of this
report. 
Below we discuss ways of adding new modules or replacing modules which
already exist in the code.

To close the system of transport equations \req{main} we need, in
addition to the sources and sinks of \eqs\ref{S})--(\ref{P}), the
transport matrix of \req{fluxes}, the electrical conductivity
$\sigma_\pll$ and the driven current $j_{CD}$, specify also the surface
functions $V',I,G_1,G_2,G_3$.
Then with appropriate initial conditions and boundary conditions, we can
determine the time evolution of the radial distributions of the electron
density $n_e$, temperatures $T_{e,i}$ and current density.

\subsection{Initial conditions}

Initial conditions do not usually play a significant role in tokamak
transport modeling, because normally steady-state discharge conditions
are under study. 
However, for the analysis of transient processes, such as plasma current
ramp-up, initial conditions may play an essential role. 
When available, measured experimental quantities may be used to supply
distributions, but these are not always available, especially for the
poloidal flux $\psi$ and the corresponding current density. 
The simulation of some processes can be sensitive to the initial choice
of the current density profile, and this choice may require care.
It is usually sufficient to set initial conditions self-consistently,
to give smooth plasma evolution during the initial moments of the run. 

For the first three equations in \req{main} the natural initial
conditions are 
\begin{equation}\label{ic}
\begin{array}{l}
\lb.n_e(\rho,t)\rb|_{t=0}=n_{e0}(\rho)
,\bigskip\\
\lb.T_e(\rho,t)\rb|_{t=0}=T_{e0}(\rho)
,\bigskip\\
\lb.T_i(\rho,t)\rb|_{t=0}=T_{i0}(\rho)
.\end{array}
\end{equation}
A similar initial condition for the poloidal flux would read
$\lb.\psi(\rho,t)\rb|_{t=0}=\psi_0(\rho)$. 
However, this is of no practical use because it is usually more
convenient to prescribe the current density $j_\pll$ or the rotational
transform $\mu,$ rather than the poloidal flux $\psi$ itself. 
Thus, appropriate initial conditions are alternative 
\begin{equation}\label{icj0}
\begin{array}{l}
{\rm either}\qquad
\lb.j_\pll(\rho,t)\rb|_{t=0}=j_0(\rho)
\qquad\mbox{or}\qquad
\lb.\mu(\rho,t)\rb|_{t=0}=\bar\mu(\rho)
.\end{array}
\end{equation}
These conditions involve either the second or the first derivative of
the unknown function $\psi(\rho)$, respectively. 
With either initial condition, $\psi(\rho)$ must be determined using
either \req{dfjp} or \req{dfmu}.
A further well-defined quantity is the total plasma current $I_{\rm pl},$
which is always measured in experiments, while the current distribution
during the early phase of a tokamak discharge is usually not known.
Therefore it is reasonable to require that $I_{\rm pl}$ takes a given
value and the normalizations of the current density or rotational
transform profiles are adjusted to be self--consistent using
\begin{equation}\label{icj}
\lb\{\begin{array}{l}\dstyle
\lb.j_\pll(\rho,t)\rb|_{t=0}=j_0(\rho)
\bigskip\\ \dstyle
\int\limits_0^{\rho_B}J^{-2}j_0 V'd\rho=2\pi R_0 I_{\rm pl}
\end{array}\rb.
\end{equation}
or
\begin{equation}\label{icmu}
\lb.\mu(\rho,t)\rb|_{t=0}=\frac{\bar\mu(\rho)}{\bar\mu(\rho_B)}
\frac{\mu_0 I_{\rm pl}}{2\pi B_0 \rho_B G_2(\rho_B)}
,\end{equation}
respectively.
The boundary condition $J(\rho_B)=1$ is used in \req{icj}.

Also convenient and useful is the initial condition with the plasma
current distributed according to the steady-state condition 
$\prtii{\psi}{t}{\rho}=0\Rightarrow\prti{\psi}{t}
=U_{\rm pl}(\rho)=Const$. 
Using the parallel Ohm's law \req{Ohm} and assuming
$\dot B_0(t=0)=0$ we re-write the condition as 
\begin{equation}\label{icU}
\lb.\Bigl(j_\pll-j_{BS}-j_{CD}\Bigr)\rb|_{t=0}
=\frac{2\pi\rho}{V'}\sigma_\pll(\rho)U_{\rm pl}
=C_\sigma\frac{\rho}{V'}\sigma_\pll(\rho)
.\end{equation}
Similar to \eqs\ref{icj}),~(\ref{icmu}) the factor $C_\sigma$ in
\req{icU} should be found from the condition that the total current is
$I_{\rm pl}$. 
Although the current relaxation time is long in a tokamak, the
stationary initial condition \req{icU} is widely used when the direct
measurements of the current density are not available. 
The condition of \req{icU} is also physically relevant when
non-diffusive processes cause fast current relaxation, or when the 
preceding long phase of relaxation is of no interest. 

\subsection{Boundary conditions for densities and temperatures}

The boundary conditions for \req{main} at the magnetic axis $\rho=0$ are
imposed by the geometry of problem and the requirement that all fluxes
should vanish at $\rho=0$ 
\begin{equation}\label{bc0}
\lb.\Gamma_e\rb|_{\rho=0}
=\lb.q_e\rb|_{\rho=0}
=\lb.q_i\rb|_{\rho=0}
=0
.\end{equation}

Boundary conditions at $\rho=\rho_B$ for the first three equations of
the system \req{main} usually take one of two forms 
\begin{equation}\label{bc}
\begin{array}{ll}
n_e(\rho_B)=n_{eB}(t)
\qquad &\mbox{or}\Quad
\Gamma_e(\rho_B)=\Gamma_{eB}(t)
,\bigskip\\
T_e(\rho_B)=T_{eB}(t)
\qquad &\mbox{or}\Quad
q_e(\rho_B)=q_{eB}(t)
,\bigskip\\
T_i(\rho_B)=T_{iB}(t)
\qquad &\mbox{or}\Quad
q_i(\rho_B)=q_{iB}(t)
.\end{array}
\end{equation}
All the functions on the right hand sides in \req{bc}, besides
explicitly depending on time $t,$ can also depend on all other plasma
parameters. 
The boundary conditions of \req{bc} can be imposed in any combination as
described in Section~4.4.

It worth noting now that, in general, the question of boundary
conditions is not quite straightforward. 
Indeed, in the configuration space, the plasma boundary is supposed to
coincide with some flux surface $\mbox{\msym S}_B.$ 
This boundary surface serves as an input to a solver of the
Grad-Shafranov equilibrium equation. 
The solver returns a position of this surface in flux coordinate $\rho_B.$
Even in the simplest case, when the plasma boundary
$\mbox{\msym S}_B$ is fixed in space and does not move, the
corresponding value $\rho_B$ can vary in time depending on the plasma
pressure and current density distributions. 
Therefore, the plasma boundary $\rho_B(t)$ moves in the frame of the
Lagrangian flux coordinate $\rho.$ 
Thus the problem under consideration is always a moving boundary
problem and this is discussed in more detail in Section~3.11.

It is also possible that the required boundary conditions cannot be
expressed explicitly in any of forms given in \req{bc}. 
This can be the case when more elaborated boundary conditions are
imposed by considering SOL models.
Correct description of these cases could require the use of MHD or
kinetic theory. 
In the Astra code, these situations are implemented by calling special
subroutines as described in Section 4. 

\subsection{Boundary condition for the poloidal flux}

The left boundary condition for the fourth equation in \req{main} is
straightforward, 
${\lb.\prti{\psi}{\rho}\rb|_{\rho=0}\hspace*{-10pt}=0}.$
The boundary condition at $\rho=\rho_B$ is more complicated. 
Generally, the condition should be obtained by solving a magnetostatic
problem in the exterior region (with respect to the plasma) taking into
account all poloidal currents with their mutual inductances,
magnetization of the iron core, currents induced in a vacuum chamber and
so on. 
However, this would unreasonably complicate the problems under
consideration. 
It is more practical to use one of the three possibilities described
below. 

\subsubsection{Prescribed plasma current} 
This is the most frequently used boundary condition for transport
modeling. 
%The condition is justified by a relative independence of the equation
%for $\psi$ from other transport equations.  
%Since the processes of particle and energy transport have the
%characteristic times much shorter than this for the poloidal flux they
%can be studied assuming prescribed or even time-independent plasma
%current.  
Making use of \req{Ipl} we write the condition as
\begin{equation}\label{bcI}
\lb.\prti{\psi}{\rho}\rb|_{\rho=\rho_B}
=\frac{\mu_0}{G_2(\rho_B)}I_{\rm pl}(t)
\end{equation}
where $I_{\rm pl}(t)$ is the total plasma current with a prescribed time
dependence. 

\subsubsection{Prescribed loop voltage} 
This is usually not a good choice of boundary condition. 
For instance, in Ohmic plasmas this boundary condition gives rise to a
thermal instability. 
On the other hand, this instability is very slow and can be easily
stabilized. 
Therefore, this condition is also included in the code as a possible
option. 
In the simplest case where the plasma boundary is fixed (see discussion
in Section~3.4) it reads 
\begin{equation}\label{bcU}
\lb.\prti{\psi}{t}\rb|_{\rho=\rho_B}
=U_{\rm plB}(t)
\end{equation}
with $U_{\rm plB}$ being the prescribed boundary loop voltage.

\subsubsection{External circuit equation} 
This provides a boundary condition which includes the two previous
conditions as particular cases. 
It uses a simplified description of an external circuit with the single
equation 
\begin{equation}\label{bcL}
U_{\rm ext}=\frac{d}{dt}\lb(L_{\rm ext}I_{\rm pl}\rb)+U_{\rm plB}
\end{equation}
where $L_{\rm ext}$ is the external inductance of the plasma, and 
$U_{\rm plB}=U_{\rm pl}(\rho=\rho_B)$
is the loop voltage calculated at the plasma edge \req{Uplb}.
$U_{\rm ext}$ is a given function of time associated with the voltage
produced by the primary coil. 

\subsection{Closure of the equilibrium and transport equations}

The four equations of \req{main} do not comprise a closed system, even
when all the fluxes and right hand sides are determined. 
Indeed, the functions $V(\rho,t)$, $J(\rho,t)$, $G_{1,2}(\rho,t)$
included in \req{main} still remain unknown. 
Moreover, the independent variable $\rho$ is also an unknown
function of space coordinates and time: $\rho=\rho(r,z,t)$. 
It is the solution of the Grad-Shafranov equation \req{GSh} which gives
the instant relations of all flux coordinates, such as $\rho$, $V$, $J$,
etc., to one another and to the laboratory coordinate system $\{r,z\}$. 
These relations can vary in time due to time evolution of plasma
parameters as described by \req{main}. 
Hence to study the plasma evolution in a tokamak we should solve the set
of one-dimensional transport equations (\ref{main}),~(\ref{fluxes})
simultaneously with the two-dimensional equilibrium equation (\ref{GSp}).
\footnote{Because of this combination of 1D and 2D equations
the transport codes are often referred to as 1.5D codes.}

First of all, we limit our consideration to the internal equilibrium
problem with a prescribed plasma boundary. 
This boundary is not supposed to be fixed in time so that
the entire plasma can move according to a given law.
However, we do not consider the external fields and currents which are
required to provide this boundary evolution and the plasma movement
because the solving of external equilibrium problem is beyond the scope
of the Astra code.  

The closing procedure can be seen as follows.
At each time step, the solution to the transport equations
(\ref{main}),~(\ref{fluxes}) provides the functions $p(\rho)$ and
$j_\pll(\rho)$ on the \rhs of \req{GSp}.
These functions are determined on the interval $0\leq\rho\leq\rho_B$.
On the other hand, the solution of \req{GSp} is sought within the
prescribed plasma boundary $\mbox{\msym S}_B.$  
However, the requirement that this boundary $\mbox{\msym S}_B$ must
coincide with the flux surface $\rho=\rho_B$ would overdetermine the
problem. 
According to the particular experimental conditions, it is possible to
adopt different ways of treating this contradiction.
For instance, the overdetermination can be removed if the surface
$\mbox{\msym S}_B$ is given by a set of points $\{r_i,z_i\}$ while in
between it is allowed to be arbitrary.
This approach is practically useless because it suppose a corrugated
plasma surface which most probably is not compatible with the external
magnetic system.

More practical would be to consider the whole set of boundary parameters 
$\{\rho_B,r_i,z_i\}$ as approximate which has to be fitted with some
accuracy. 
If $\{r_i,z_i\}$ is firmly prescribed and $\rho_B$ is free then the
formulation will be referred as  the prescribed plasma boundary. 
The case with fixed $\rho_B$ and flexible $\{r_i,z_i\}$ will be called
the adjustable plasma boundary. 
These two approaches are discussed below. 

\subsubsection{Prescribed plasma boundary}

Presently, the following procedure is implemented in the Astra code.
The equilibrium equation is solved requiring that
$\mbox{\msym S}_B$ coincides with one of the $\rho$-surfaces, say
$\rho=\hat\rho_B$, which can be either external or internal with respect
to $\rho_B$. 
In the case $\rho_B<\hat\rho_B$, there is a gap between the current
carrying plasma and the presumed boundary $\mbox{\msym S}_B$. 
In the case $\rho_B>\hat\rho_B$ the excessive plasma layer should be
scraped off. 
In both cases, the functions $p(\rho)$ and $j_\pll(\rho)$ are re-defined
on the interval between $\rho_B$ and $\hat\rho_B$ so that in the region
$\min(\rho_B,\hat\rho_B)\leq\rho\leq\max(\rho_B,\hat\rho_B)$:
$p(\rho)=p(\rho_B)$ and $j_\pll(\rho)=0.$ 
In addition, a surface skin current is added at $\rho=\hat\rho_B$ in
order to keep the total plasma current unchanged. 
Solving \req{GSp} is then iteratively repeated until $\hat\rho_B$
remains constant with sufficient accuracy. 
The new boundary position is selected as $\rho=\hat\rho_B$. 
The described procedure does not change when a plasma boundary moves or
varies its shape. 

In practical calculations, at each time step it is sufficient to use the
value of $\hat\rho_B$ returned by the equilibrium solver after one
iteration as a new position of the boundary point $\rho_B$. 
Indeed, as it has already been discussed, modifications in current
density, pressure and the corresponding changes in plasma
diamagnetism cause time variation of the plasma size. 
Consequently, the time variation of $\rho_B(t)$ is of the order
$\dot\beta_{tor}\Delta t$ and hence is very small. 
Typically, the relative variation of $\rho_B$ is as small as $10^{-3}$
per time step. 

\subsubsection{Adjustable boundary}

More relevant would be to consider the prescribed boundary shape
as approximate which has to be fitted with some accuracy. 
Then the different points on the boundary can be matched with different
admissible errors according to their confidence weight. 
%(see Section~3.11.2). 
A trivial example of such an approach would be to fix the innermost point
on the boundary (associated with internal limiter) and, adjusting the
rest of the boundary $\mbox{\msym S}_B,$ allow the entire plasma to
shift in order to avoid a contact with a limiter and preserve $\rho_B$
unchanged. 

This option assumes that the plasma boundary is given with one free
parameter, which may be the boundary elongation, plasma shift, etc. 
When plasma is expanding then the freedom is used in order to fulfill the
condition that the plasma remains within the flux surface $\rho_B$ and
thus avoids plasma contact with a limiter. 
All other features of the code remain the same as in the case with
prescribed boundary shape. 
This version is easier for a numerical implementation and, probably,
more relevant from the physical point of view. 
However, it requires a specification of the admissible freedom. 
We note in conclusion that all the details have significance which is
more formal than practical, because the real variation in $\rho_B(t)$ is
negligibly small. 

\subsection{Auxiliary transport equations}

Usually, a tokamak plasma contains a variety of different species. 
These might be hydrogen isotopes, helium ash, or other impurities. 
Each of them can be characterized by its own density and temperature. 
Non-Maxwellian populations of the main plasma components can also be
treated as independent species. 
Densities, temperatures, velocities of rotation and other thermodynamic
characteristics of plasmas are believed to obey the law of collisional
diffusion. 
For treatment of all these processes not included in \req{main} the
Astra code provides a set of additional transport equations. 
The present practice of transport tokamak modeling never uses the
transport matrix of \req{fluxes} in full and most present-day tokamaks
work with a time-independent toroidal magnetic field. 
Therefore, these additional equations have a truncated form 
\begin{equation}\label{aux}
%\hspace*{-5mm}
\prti{}{t}\left(V'f_j\right)
=\prti{}{\rho}
\left[V'
\av{(\nb\rho)^2}\left(D_j\prti{f_j}{\rho}-v_j f_j\right)
\right]
+V'S_j
\qquad\qquad 
j=1,2,3
.\end{equation}
In version 5.0 of the Astra code the number of equations (\ref{aux})
cannot exceed 3, but this restriction can easily be lifted. 
The initial and boundary conditions for \req{aux} are similar to those
for the first three equations in \req{main} 
\begin{equation}\label{icf}
\lb.f_j(\rho,t)\rb|_{t=0}=f_{j0}(\rho)
\end{equation}
and
\begin{equation}\label{bcf}
\lb\{\begin{array}{l}
\lb.\left(D_j\prti{f_j}{\rho}-v_j f_j\right)\rb|_{\rho=0}
=0
,\bigskip\\
f_j(\rho_B,t)=f_{jB}(t)
\qquad \mbox{or}\Quad
\lb.\left(-D_j\prti{f_j}{\rho}+v_j f_j\right)\rb|_{\rho_B}=\Gamma_{jB}(t)
.\end{array}\rb.
\end{equation}

\subsection{Equation for gas puff neutrals} 

The subroutine solves a kinetic equation for a neutral distribution
function $f_N$ 
\begin{equation}\label{fN}
%\begin{array}{l}\dstyle
%{\bf v}\nb f_N+\lb(s_{ion}+s_{cx}\rb)f_N=\lb(s_{cx}N+s_{rec}n_e\rb)f_{Mi}
v\prti{f_N}{x}+\lb(s_{ion}^{(e)}n_e+s_{ion}^{(i)}n_i+s_{cx}n_i\rb)f_N
%\bigskip\\ \hspace*{10em}\dstyle
=\frac{\sqrt{3}}{2}n_i\lb(s_{cx}N+s_{rec}n_e\rb)
\delta\lb(v\pm\frac{v_{Ti}}{\sqrt{3}}\rb)
,%.\end{array}
\end{equation}
where $v_{Ti}=\sqrt{2T_i/m_i}.$
The factor $\sqrt{3}$ is included in \req{fN} in order to describe
isotropic velocity distribution of the charge exchange neutrals.
The equation (\ref{fN}) is solved in a slab geometry assuming that the
plasma slab thickness is equal to twice the  minor radius in the
equatorial cross-section $2a_B$. 
This approximation is reasonable when the mean free path of a neutral is
much smaller than $a_B$. 
In present-day tokamaks, the condition is usually fulfilled with a large
margin. 
As long as the transit time of the neutral atoms is quite small the
problem is treated as steady state.

It is assumed that the incoming neutral particles have a
velocity distribution 
\begin{equation}\label{fN0}
\begin{array}{rl}\dstyle
\Bigl.f_N(x=-a_B,v)\Bigr|_{v>0}
&\dstyle
=N_1\delta\lb(v-v_1\rb)
+N_2\delta\lb(v-v_2\rb)
,\bigskip\\ \dstyle
\Bigl.f_N(x=a_B,v)\Bigr|_{v<0}
&\dstyle
=N_1\delta\lb(v+v_1\rb)
+N_2\delta\lb(v+v_2\rb)
,\end{array}
\end{equation}
with $v_{1,2}=\sqrt{2E_{1,2}/m_i}.$
%where $f_{Mi}$ is the Maxwellian distribution of ions. 
%where the factor $\sqrt{3}$ on the \rhs appears from th
\req{fN} is solved as described in \cite{DK} and the neutral density and
temperature are then found as 
\begin{equation}\label{N}
%N({\bf r})=\int f_N({\bf r},{\bf v})d^3v
N(x)=\int f_N(x,v)dv
\qquad\mbox{and}\qquad
%T_N({\bf r})=\int \frac{m_iv^2}{2}f_N({\bf r},{\bf v})d^3v
T_N(x)=\frac{1}{N(x)}\int \frac{m_iv^2}{2}f_N(x,v)dv
.\end{equation}

\subsection{Other Astra compatible packages available by request}

A number of other modules are incorporated in the Astra code.
Some of them are quite complicated and time consuming. 
The code is continually supplemented with newer ones.
Only few of these additional packages are included in the 
standard Astra version.
Sometimes the interfaces only are provided, while the packages 
can be obtained directly from their authors. 
A description of all these modules is out of the scope of this report. 
Only a list of the most useful modules with the names of the
contact persons is presented here.

\subsubsection{Additional heating and current drive}

\paragraph{Neutral beam heating and current drive \rm\cite{Pole}.} 
The NBI Heating and CD package of the Astra code was developed by
A.~R.~Polevoi (polevoy@nfi.kiae.su). 
A simplified time independent version of the package is distributed with
the Astra code. 
It is briefly described in Section~4.6.2.
A more advanced time dependent version is available. 

\paragraph{Electron cyclotron heating and current drive \rm\cite{Poli}.} 
The wave equation for the electromagnetic waves in a cold plasma is
solved in the Gaussian beam approximation. 
The wave absorption at the frequency of the electron cyclotron resonance
(ECR) and its harmonics is calculated for a Maxwellian plasma in the
weakly relativistic approximation. 
The driven current is estimated using the adjoint technique.
This package
%, called \tws{TORBEAM}, 
is developed by E.~Poli 
(Emanuele.Poli@ipp.mpg.de, see also 
http://www.aug.ipp.mpg.de/\symbol{126}emp/torbeam.html).

\paragraph{Lower hybrid heating and current drive \rm\cite{EP}.} 
The wave equation is solved for the lower-hybrid frequency
range in the cold plasma approximation using the geometric optics
technique. 
The Fokker--Planck equation is solved in the 1D approximation.
Ion and $\alpha-$particle absorption is also included.
(A.~N.~Saveliev, saveliev@ans.ioffe.rssi.ru).

\paragraph{Ion cyclotron heating \rm\cite{MB}.} 
The wave equation for the IC frequency range is solved in the eikonal
approximation. 
The code takes into account minority and ion cyclotron harmonic heating,
ion-ion hybrid resonance heating and wave absorption due to coupling to
the acoustic and ion-Bernstein waves. 
The code also includes a module for the antenna-plasma coupling which
provides the boundary conditions for the ray tracing. 
(M.~Brambilla, Marco.Brambilla@ipp.mpg.de).

\subsubsection{Transport coefficients}

The Astra library includes about one hundred different neoclassical
and anomalous transport coefficients which can be expressed as formulae. 
Recently, more complicated routines for evaluation of the transport
coefficients were developed. 
Some of them can be included in the Astra code as separate modules. 

\paragraph{Neoclassical transport {\tt NCLASS} \rm\cite{NCLASS}.}
A matrix of neoclassical transport coefficients for a multi-species
tokamak plasma is calculated.
(W.A. Houlberg, houlbergwa@ornl.gov)

\paragraph{Turbulent transport}
The following three transport models provide transport matrix due to
toroidal ion temperature gradient (ITG) modes and trapped electron modes
(TEM). 

-- { Weiland-Nordman model \rm\cite{WN}.}  
(J.~Weiland, elfjw@elmagn.chalmers.se). 

-- { IFS/PPPL model \rm\cite{ifspppl}.} (W.~Dorland,
bdorland@pppl.gov). 

-- { GLF23 model \rm\cite{glf}.} (J.~Kinsey,
kinsey@apollo.gat.com).

%\paragraph{Itoh-Fukuyama model \rm\cite{fitoh}.}  (D.~Mikkelsen, mikk@pppl.gov). 

\subsubsection{Impurity transport and radiation}
The subroutine {\tt STRAHL} \cite{strahl} describing transport and
radiation of impurity ions (R.~Dux, Dux@ipp.mpg.de,
http://www.aug.ipp.mpg.de/\symbol{126}Ralph.Dux/home.html) 
can be linked with the code Astra. 
No cross influence of impurities on the main plasma transport is
included. 

%\subsubsection{Scrape-off layer description}

\subsubsection{Stability analysis}

The code provides interfaces to stability codes as {\tt PEST,
MAST, CASTOR, GARBO}. 
For this codes, a special input file is written for specified time
slices which can be used for a post-run stability analysis.
In addition, more simplified routines are available for on line
evaluation of instability induced transport. 

\paragraph{Sawtooth oscillations ({\rm subroutine \tt MIXINT}.)} 
Sawtooth description \cite{PP} based on the Ka\-domtsev theory for
magnetic field line reconnection \cite{BB,PP} is included in the standard
Astra distribution. 

\paragraph{Tearing modes ({\rm subroutine \tt ISLAND}.)} 
Simplified $\Delta'$ analysis for saturated island width in a
cylindrical plasma is provided (A.N.Chudnovskij, chdn@wowa.net.kiae.su). 

%\hyphenation
%\paragraph{Balloonning modes.} (H.-P.Zerhfeld, zehrfeld@ipp.mpg.de)

\vfill\eject
